\documentclass[a4paper,12pt]{article}

% -------------------------------------------------
% Кодировка и язык
% -------------------------------------------------
\usepackage[T2A]{fontenc}
\usepackage[utf8]{inputenc}
\usepackage[russian]{babel}

% -------------------------------------------------
% Математика
% -------------------------------------------------
\usepackage{amsmath,amssymb,amsthm,mathtools}
\usepackage{icomma}

% -------------------------------------------------
% Шрифт и типографика
% -------------------------------------------------
\usepackage{helvet}
\renewcommand{\familydefault}{\sfdefault}
\usepackage{microtype}
\usepackage{setspace}
\onehalfspacing
\usepackage{parskip}

% -------------------------------------------------
% Геометрия страницы
% -------------------------------------------------
\usepackage[
    a4paper,
    left=20mm,
    right=20mm,
    top=20mm,
    bottom=22mm
]{geometry}

% -------------------------------------------------
% Графика
% -------------------------------------------------
\usepackage{tikz}
\usetikzlibrary{calc}
\usepackage{tkz-euclide}
\usepackage{pgfplots}
\pgfplotsset{compat=1.18}

% -------------------------------------------------
% Таблицы и списки
% -------------------------------------------------
\usepackage{tabularx,booktabs,longtable}
\usepackage{enumitem}

% -------------------------------------------------
% Цвета
% -------------------------------------------------
\usepackage{xcolor}

\definecolor{accent}{HTML}{008080}
\definecolor{darkaccent}{HTML}{004D4D}
\definecolor{textgray}{HTML}{333333}
\definecolor{softgray}{HTML}{777777}
\definecolor{lightaccent}{HTML}{E8F2F2}

% -------------------------------------------------
% Ссылки
% -------------------------------------------------
\usepackage{hyperref}
\hypersetup{
    colorlinks=true,
    linkcolor=darkaccent,
    urlcolor=darkaccent,
    citecolor=darkaccent
}

% -------------------------------------------------
% Колонтитулы
% -------------------------------------------------
\usepackage{fancyhdr}
\setlength{\headheight}{15pt}

\pagestyle{fancy}
\fancyhf{}
\fancyhead[L]{\small\textcolor{softgray}{Числа и периодические дроби}}
\fancyhead[R]{\small\textcolor{softgray}{10 класс}}
\fancyfoot[C]{\small\textcolor{softgray}{\thepage}}

\renewcommand{\headrulewidth}{0.4pt}
\renewcommand{\headrule}{
    \hbox to\headwidth{
        \color{accent}\leaders\hrule height \headrulewidth\hfill
    }
}

% -------------------------------------------------
% Настройки списков
% -------------------------------------------------
\setlist[itemize]{
    leftmargin=6mm,
    itemsep=3pt,
    topsep=4pt
}

\setlist[enumerate]{
    leftmargin=8mm,
    itemsep=6pt,
    topsep=5pt
}

% -------------------------------------------------
% Служебные команды
% -------------------------------------------------
\newcommand{\checkitem}{\(\square\)}
\newcommand{\N}{\mathbb{N}}
\newcommand{\Z}{\mathbb{Z}}
\newcommand{\Q}{\mathbb{Q}}
\newcommand{\R}{\mathbb{R}}

\newcommand{\key}[1]{%
    \par\medskip
    {\color{darkaccent}\bfseries #1}
    \par\smallskip
}

\newcommand{\ruleline}{%
    \par\medskip
    {\color{accent}\hrule height 0.7pt}
    \medskip
}

\begin{document}

% =================================================
% ТИТУЛЬНЫЙ ЛИСТ
% =================================================

\begin{titlepage}
\thispagestyle{empty}

\begin{center}

\vspace*{23mm}

{\color{accent}\rule{0.26\linewidth}{1.2pt}}

\vspace{12mm}

{\fontsize{29}{34}\selectfont\bfseries
Чек-лист
}

\vspace{6mm}

{\fontsize{20}{25}\selectfont\bfseries
Рациональные и иррациональные числа
}

\vspace{3mm}

{\fontsize{17}{22}\selectfont
Периодические десятичные дроби
}

\vspace{15mm}

{\large
10 класс
}

\vfill

{\large\bfseries
Лёвин Артём Александрович
}

\vspace{2mm}

{\large
эксклюзивно для Ани Трапезниковой
}

\vspace{9mm}

{\large
\today
}

\vspace{25mm}

\end{center}

% Сдержанная кривая Безье внизу титульного листа
\begin{tikzpicture}[remember picture,overlay]
    \draw[
        color=accent,
        line width=1.4pt
    ]
    ([xshift=17mm,yshift=15mm]current page.south west)
    .. controls
    ([xshift=58mm,yshift=31mm]current page.south west)
    and
    ([xshift=-66mm,yshift=4mm]current page.south east)
    ..
    ([xshift=-17mm,yshift=18mm]current page.south east);
\end{tikzpicture}

\end{titlepage}

% =================================================
% ОСНОВНОЙ ЧЕК-ЛИСТ
% =================================================

\section*{\textcolor{darkaccent}{1. Что нужно знать о десятичных дробях}}

\ruleline

\key{Конечная десятичная дробь}

Это десятичная дробь, у которой после запятой записано конечное
число цифр.

Например,
\[
0{,}6,\qquad
-0{,}75,\qquad
3{,}125.
\]

Любая конечная десятичная дробь является рациональным числом.

\key{Бесконечная периодическая десятичная дробь}

Если одна цифра или группа цифр повторяется бесконечно много раз,
повторяющуюся часть называют \textbf{периодом}.

Например,
\[
0{,}6666\ldots=0{,}(6),
\]
\[
0{,}727272\ldots=0{,}(72),
\]
\[
0{,}285714285714\ldots=0{,}(285714).
\]

\begin{itemize}
    \item[\checkitem] Скобки охватывают весь повторяющийся фрагмент.
    \item[\checkitem] В периоде может находиться одна или несколько цифр.
    \item[\checkitem] Любая периодическая десятичная дробь является рациональным числом.
\end{itemize}

\subsection*{\textcolor{darkaccent}{Как получить десятичную запись обыкновенной дроби}}

Числитель делят на знаменатель.

\[
\frac{2}{3}=0{,}(6),
\qquad
\frac{8}{11}=0{,}(72),
\]
\[
\frac{3}{5}=0{,}6,
\qquad
-\frac34=-0{,}75.
\]

Деление может закончиться либо привести к повторению остатка.
Повторение остатка означает появление периода.

% -------------------------------------------------

\section*{\textcolor{darkaccent}{2. Из периодической дроби в обыкновенную}}

\ruleline

Существуют два основных случая.

\subsection*{\textcolor{darkaccent}{Случай 1. Период начинается сразу после запятой}}

Пусть
\[
x=0{,}(\underbrace{a_1a_2\ldots a_k}_{k\text{ цифр}}).
\]

В числитель записываем число, составленное из цифр периода.

В знаменатель записываем столько девяток, сколько цифр содержится
в периоде:

\[
0{,}(A)=\frac{A}{10^k-1}.
\]

Например,
\[
0{,}(15)=\frac{15}{99}=\frac{5}{33},
\]
\[
0{,}(27)=\frac{27}{99}=\frac{3}{11},
\]
\[
0{,}(123)=\frac{123}{999}=\frac{41}{333}.
\]

\begin{itemize}
    \item[\checkitem] Количество девяток равно количеству цифр периода.
    \item[\checkitem] После получения дроби обязательно проверяйте возможность сокращения.
\end{itemize}

% -------------------------------------------------

\subsection*{\textcolor{darkaccent}{Случай 2. Перед периодом есть цифры}}

Рассмотрим
\[
x=0{,}23(18).
\]

После запятой имеются:

\[
\underbrace{23}_{\text{до периода}}
\underbrace{18}_{\text{период}}.
\]

В числителе записываем
\[
2318-23.
\]

В знаменателе:

\begin{itemize}
    \item две девятки --- по числу цифр периода;
    \item два нуля --- по числу цифр между запятой и периодом.
\end{itemize}

Получаем
\[
0{,}23(18)
=
\frac{2318-23}{9900}
=
\frac{2295}{9900}
=
\frac{51}{220}.
\]

\key{Общее правило}

Если имеется дробь
\[
0{,}a_1a_2\ldots a_m(b_1b_2\ldots b_k),
\]
то
\[
x=
\frac{B-A}
{10^m(10^k-1)},
\]
где

\[
A=\text{число из цифр до периода},
\]
а
\[
B=\text{число из цифр до периода и первой записи периода}.
\]

В знаменателе получается:

\[
\underbrace{99\ldots9}_{k\text{ девяток}}
\underbrace{00\ldots0}_{m\text{ нулей}}.
\]

\subsection*{\textcolor{darkaccent}{Почему правило работает}}

Для
\[
x=0{,}23(18)
\]
имеем
\[
100x=23{,}(18),
\]
\[
10000x=2318{,}(18).
\]

Вычитаем первое равенство из второго:

\[
10000x-100x=2318-23,
\]
\[
9900x=2295,
\]
откуда
\[
x=\frac{2295}{9900}
=\frac{51}{220}.
\]

Такой способ особенно полезен, если требуется строго обосновать
формулу.

% -------------------------------------------------

\section*{\textcolor{darkaccent}{3. Рациональные числа}}

\ruleline

\key{Определение}

Число называется \textbf{рациональным}, если его можно представить
в виде
\[
\frac{m}{n},
\qquad
m\in\Z,\quad n\in\Z,\quad n\ne0.
\]

Множество рациональных чисел обозначают
\[
\Q.
\]

К рациональным относятся:

\[
5=\frac51,
\]
\[
-\frac37,
\]
\[
0{,}6=\frac35,
\]
\[
0{,}(6)=\frac23.
\]

\begin{itemize}
    \item[\checkitem] Целые числа рациональны.
    \item[\checkitem] Обыкновенные дроби рациональны при ненулевом знаменателе.
    \item[\checkitem] Конечные десятичные дроби рациональны.
    \item[\checkitem] Бесконечные периодические десятичные дроби рациональны.
\end{itemize}

\key{Критерий по десятичной записи}

Рациональное число имеет десятичную запись, которая либо конечна,
либо бесконечна и периодична.

% -------------------------------------------------

\section*{\textcolor{darkaccent}{4. Иррациональные числа}}

\ruleline

\key{Определение}

\textbf{Иррациональное число} --- действительное число, которое нельзя
представить в виде
\[
\frac{m}{n},
\qquad m,n\in\Z,\quad n\ne0.
\]

Его десятичная запись бесконечна и непериодична.

Например,
\[
\sqrt2
=
1{,}41421356\ldots
\]
является иррациональным числом.

Также иррациональны, например,
\[
\sqrt3,\qquad
\sqrt5,\qquad
\pi.
\]

Но наличие корня ещё не гарантирует иррациональность:

\[
\sqrt4=2,
\qquad
\sqrt{81}=9.
\]

Оба результата рациональны.

\key{Главный приём}

Перед тем как определять вид числа, сначала упростите выражение.

Например,
\[
(\sqrt8-3)(\sqrt8+3).
\]

Так как
\[
\sqrt8=2\sqrt2,
\]
используем формулу
\[
(a-b)(a+b)=a^2-b^2:
\]

\[
(\sqrt8-3)(\sqrt8+3)
=
(\sqrt8)^2-3^2
=
8-9=-1.
\]

Следовательно, исходное выражение имеет рациональное значение.

% -------------------------------------------------

\subsection*{\textcolor{darkaccent}{Ещё один типичный пример}}

Рассмотрим
\[
(\sqrt3-1)^2+2\sqrt3.
\]

Используем формулу
\[
(a-b)^2=a^2-2ab+b^2.
\]

Тогда
\[
(\sqrt3-1)^2
=
3-2\sqrt3+1
=
4-2\sqrt3.
\]

Следовательно,
\[
(\sqrt3-1)^2+2\sqrt3
=
4.
\]

Корни полностью сократились, поэтому значение рационально.

\key{Типичная ошибка}

В выражении
\[
(\sqrt3-1)^2
\]
роль \(b\) играет число \(1\).

Поэтому
\[
(\sqrt3-1)^2
=
3-2\sqrt3+1,
\]
а знак среднего слагаемого остаётся отрицательным.

% -------------------------------------------------

\section*{\textcolor{darkaccent}{5. Множества чисел}}

\ruleline

Основные множества связаны включениями

\[
\boxed{
\N\subset\Z\subset\Q\subset\R
}
\]

В школьном курсе обычно принимают
\[
\N=\{1,2,3,\ldots\}.
\]

Поэтому
\[
0\notin\N.
\]

Целые числа:
\[
\Z=
\{\ldots,-3,-2,-1,0,1,2,3,\ldots\}.
\]

Рациональные числа:
\[
\Q=
\left\{
\frac{m}{n}\;:\;
m,n\in\Z,\ n\ne0
\right\}.
\]

Действительные числа:
\[
\R.
\]

Множество иррациональных чисел можно записать как
\[
\R\setminus\Q.
\]

\begin{center}
\begin{tikzpicture}[scale=0.93]
    \draw[accent,line width=1pt]
        (0,0) ellipse (1.3cm and 0.8cm);
    \draw[accent,line width=1pt]
        (0,0) ellipse (2.4cm and 1.5cm);
    \draw[accent,line width=1pt]
        (0,0) ellipse (3.6cm and 2.2cm);
    \draw[accent,line width=1pt]
        (0,0) ellipse (4.8cm and 2.9cm);

    \node at (0,0) {$\N$};
    \node at (0,1.12) {$\Z$};
    \node at (0,1.84) {$\Q$};
    \node at (0,2.54) {$\R$};

    \node[font=\small] at (0,-0.4) {$1,\ 2,\ 3,\ldots$};
    \node[font=\small] at (-1.7,-0.95) {$0,\ -4$};
    \node[font=\small] at (2.35,-1.55) {$\dfrac37,\ 0{,}(6)$};
    \node[font=\small] at (-3.45,-2.2) {$\sqrt2,\ \sqrt3$};
\end{tikzpicture}
\end{center}

\key{Как определять множество числа}

Удобно указывать \textbf{наименьшее} из стандартных множеств,
к которому принадлежит число.

Например,
\[
7\in\N,
\]
хотя одновременно
\[
7\in\Z,\qquad 7\in\Q,\qquad 7\in\R.
\]

Для числа \(-5\) наименьшее подходящее множество:
\[
-5\in\Z.
\]

Для
\[
\frac27
\]
это множество
\[
\Q.
\]

Для
\[
\sqrt2
\]
это
\[
\R\setminus\Q.
\]

% -------------------------------------------------

\section*{\textcolor{darkaccent}{6. Что особенно важно запомнить}}

\ruleline

\begin{itemize}
    \item[\checkitem]
    При делении обыкновенной дроби на десятичную запись следите за
    повторением остатков.

    \item[\checkitem]
    В записи \(0{,}(72)\) периодом является сразу число \(72\).

    \item[\checkitem]
    Для \(0{,}(A)\) в знаменателе дроби появляются только девятки.

    \item[\checkitem]
    Для смешанной периодической дроби в знаменателе сначала идут
    девятки по числу цифр периода, затем нули по числу цифр до периода.

    \item[\checkitem]
    Полученную обыкновенную дробь нужно сократить.

    \item[\checkitem]
    Знаменатель обыкновенной дроби всегда должен быть ненулевым.

    \item[\checkitem]
    Рациональная десятичная запись либо конечна, либо периодична.

    \item[\checkitem]
    Иррациональная десятичная запись бесконечна и непериодична.

    \item[\checkitem]
    Корень следует сначала упростить: \(\sqrt{8}=2\sqrt2\),
    \(\sqrt{18}=3\sqrt2\).

    \item[\checkitem]
    Перед классификацией сложного выражения сначала выполните
    допустимые преобразования.

    \item[\checkitem]
    Помните:
    \[
    \N\subset\Z\subset\Q\subset\R.
    \]
\end{itemize}

% =================================================
% ТРЕНИРОВОЧНЫЙ БЛОК
% =================================================

\newpage

\section*{\textcolor{darkaccent}{Тренировочный блок}}

\ruleline

\textbf{Путь А.}
Выполните 10 упражнений по порядку.
Сложность постепенно возрастает.

\subsection*{\textcolor{darkaccent}{Средний уровень}}

\begin{enumerate}

\item
Представьте в виде десятичной дроби:

\begin{enumerate}[label=\arabic*)]
    \item \(\dfrac23\);
    \item \(\dfrac8{11}\);
    \item \(\dfrac35\);
    \item \(-\dfrac34\).
\end{enumerate}

\item
Представьте в виде бесконечной периодической десятичной дроби:

\begin{enumerate}[label=\arabic*)]
    \item \(\dfrac5{18}\);
    \item \(\dfrac7{12}\).
\end{enumerate}

\item
Представьте в виде несократимой обыкновенной дроби:

\begin{enumerate}[label=\arabic*)]
    \item \(0{,}(15)\);
    \item \(0{,}(123)\);
    \item \(0{,}(27)\).
\end{enumerate}

\end{enumerate}

\subsection*{\textcolor{darkaccent}{Уровень выше среднего}}

\begin{enumerate}[start=4]

\item
Представьте в виде несократимой обыкновенной дроби:

\begin{enumerate}[label=\arabic*)]
    \item \(0{,}2(3)\);
    \item \(0{,}14(6)\).
\end{enumerate}

\item
Представьте в виде несократимой обыкновенной дроби:

\begin{enumerate}[label=\arabic*)]
    \item \(0{,}23(18)\);
    \item \(1{,}2(34)\).
\end{enumerate}

\item
Для каждого числа укажите наименьшее из множеств
\[
\N,\qquad \Z,\qquad \Q,\qquad \R\setminus\Q,
\]
к которому оно принадлежит:

\begin{enumerate}[label=\arabic*)]
    \item \(\sqrt{81}\);
    \item \(-12\);
    \item \(0\);
    \item \(\dfrac79\);
    \item \(\sqrt7\);
    \item \(0{,}(37)\).
\end{enumerate}

\item
Упростите выражение и определите, рационально или иррационально
его значение:

\begin{enumerate}[label=\arabic*)]
    \item \((\sqrt8-3)(\sqrt8+3)\);
    \item \((\sqrt3-1)^2+2\sqrt3\);
    \item \(\sqrt{18}-3\sqrt2\).
\end{enumerate}

\item
Упростите выражение и определите вид полученного числа:

\begin{enumerate}[label=\arabic*)]
    \item \(\sqrt{12}+\sqrt3\);
    \item \((\sqrt5+1)(\sqrt5-1)\);
    \item \(\dfrac{\sqrt{50}}{\sqrt2}\).
\end{enumerate}

\item
Определите, какие утверждения верны.
Неверные утверждения исправьте.

\begin{enumerate}[label=\arabic*)]
    \item \(0\in\N\);
    \item \(-4\in\Z\);
    \item \(\dfrac37\in\Q\);
    \item \(\sqrt2\in\R\setminus\Q\);
    \item \(\Q\subset\R\).
\end{enumerate}

\end{enumerate}

\subsection*{\textcolor{darkaccent}{Сложный уровень}}

\begin{enumerate}[start=10]

\item
Даны числа
\[
x=0{,}1(6),
\qquad
y=0{,}(3).
\]

Не используя приближённых значений корней, найдите
\[
A=(x-y)^2+2\sqrt{12}-4\sqrt3.
\]

Определите, является ли число \(A\) рациональным или иррациональным.

\end{enumerate}

% =================================================
% ОТВЕТЫ
% =================================================

\newpage

\section*{\textcolor{darkaccent}{Ответы к тренировочному блоку}}

\ruleline

\renewcommand{\arraystretch}{1.45}

\begin{table}[ht]
\centering
\caption{Ответы}
\begin{tabularx}{\linewidth}{
    >{\bfseries}p{11mm}
    X
}
\toprule
№ & Ответ \\
\midrule

1 &
1) \(0{,}(6)\);
2) \(0{,}(72)\);
3) \(0{,}6\);
4) \(-0{,}75\).
\\

2 &
1) \(0{,}2(7)\);
2) \(0{,}58(3)\).
\\

3 &
1) \(\dfrac5{33}\);
2) \(\dfrac{41}{333}\);
3) \(\dfrac3{11}\).
\\

4 &
1) \(\dfrac7{30}\);
2) \(\dfrac{11}{75}\).
\\

5 &
1) \(\dfrac{51}{220}\);
2) \(\dfrac{611}{495}\).
\\

6 &
1) \(\N\);
2) \(\Z\);
3) \(\Z\);
4) \(\Q\);
5) \(\R\setminus\Q\);
6) \(\Q\).
\\

7 &
1) \(-1\), рациональное;
2) \(4\), рациональное;
3) \(0\), рациональное.
\\

8 &
1) \(3\sqrt3\), иррациональное;
2) \(4\), рациональное;
3) \(5\), рациональное.
\\

9 &
1) неверно: \(0\notin\N\), при принятой здесь школьной договорённости;
2) верно;
3) верно;
4) верно;
5) верно.
\\

10 &
\[
x=\frac16,\qquad
y=\frac13,\qquad
A=\frac1{36}.
\]
Число \(A\) рациональное.
\\

\bottomrule
\end{tabularx}
\end{table}

\vfill

\begin{center}
\textcolor{softgray}{
\small
Перед следующим занятием отметьте задания,
в которых алгоритм ещё приходится вспоминать по образцу.
}
\end{center}

\end{document}